\documentclass[11pt,a4paper]{article}
\usepackage[margin=2.5cm]{geometry}
\usepackage[utf8]{inputenc}
\usepackage[T1]{fontenc}
\usepackage[french]{babel}
\usepackage{enumitem}
\usepackage{hyperref}

\begin{document}
	\pagestyle{empty}
	
	\noindent
	Friedly WOLI \\
	Toulouse, 31100 \\
	0767339290 \\
	\href{mailto:friedlysedjro@gmail.com}{friedlysedjro@gmail.com}
	
	\begin{flushright}
		Université Paris\-Est Créteil  \\
		61 avenue du Général de Gaulle \\
		94010 Créteil Cedex
	\end{flushright}
	
	\bigskip
	
	\textbf{Objet :} Candidature au Master Modélisation et Simulation en Mécanique des Solides \\
	
	Madame, Monsieur,\\
	
	Actuellement étudiant en mécanique énergétique à l’Université de Toulouse, je souhaite intégrer le Master « Modélisation et Simulation en Mécanique des Solides » proposé à l’Université Paris-Est Créteil. Cette formation représente pour moi une étape essentielle pour approfondir mes connaissances en modélisation numérique et en simulation des comportements complexes des matériaux et structures.
	
	Au cours de mon parcours, j’ai acquis de solides bases en mécanique des solides et des structures, en résistance des matériaux et en transferts thermiques. J’ai développé un fort intérêt pour la formulation mathématique des lois de comportement, l’analyse des contraintes et les méthodes numériques appliquées à la mécanique. Les enseignements en calcul scientifique et en programmation scientifique m’ont apporté les outils nécessaires pour aborder la modélisation avancée et la validation numérique.
	
	Le Master de l’UPEC correspond pleinement à mes aspirations académiques. Son programme, orienté vers la modélisation du comportement des matériaux, la mécanique non linéaire et la simulation numérique, constitue un cadre idéal pour développer mes compétences théoriques et scientifiques. 
	
	Mon projet professionnel est de m’orienter vers la recherche en modélisation numérique du comportement des structures et matériaux, dans la perspective d’un doctorat. Je souhaite contribuer au développement de modèles prédictifs et de méthodes de calcul performantes adaptées aux problématiques industrielles, notamment dans les secteurs aéronautique et énergétique.
	
	Rigoureux, curieux et persévérant, je suis prêt à m’investir pleinement dans cette formation exigeante, qui demande à la fois une solide maîtrise des fondements théoriques et une grande autonomie dans la mise en œuvre des outils numériques. 
	
	Je vous remercie sincèrement de l’attention que vous porterez à ma candidature et reste à votre disposition pour tout entretien.
	
	Je vous prie d’agréer, Madame, Monsieur, l’expression de mes salutations distinguées.
	
	\vspace{0.8cm}
	
	\begin{flushright}
		\textbf{Friedly WOLI}
	\end{flushright}
	
\end{document}
